\begin{abstract}
\noindent
A behaving animal that attends must run two computations at once: a \emph{policy}, deciding moment by moment whether a behaviourally relevant event has occurred, and a \emph{value estimate}, predicting how good its current situation is --- the quantity that, under any reward-driven account, supplies the teaching signal for the policy. We ask whether the way these two computations are \emph{wired together} in a recurrent attention agent, and not merely their presence, determines whether attention-like behaviour can be learned at all. We study a conv-free recurrent vision-transformer agent trained purely from sparse reward on a spatially cued orientation-change-detection task ($T=29$ frames) modelled on primate covert-attention paradigms. The agent carries two recurrent cross-attention streams: a \emph{priority stream} ($\mu$) that reads the live image and drives the actor (when to declare a change) and a \emph{value stream} ($q$) that reads accumulated memory and drives a distributional critic. We compare three wirings that differ only in pathway coupling and readout routing. With a shared recurrent memory and the priority stream commanding the actor, the agent becomes a calibrated detector (accuracy $0.868$, hit rate $0.794$, false-alarm rate $0.067$, reaction time $\sim$$1.95$ frames) and reproduces the hallmark reliability-scaled cueing signature of primate spatial attention. Severing the coupling so each stream holds a private memory abolishes learning: the agent never once declares a change (hit rate $0.000$, accuracy $0.504$, the no-change base rate). A second control that retains the coupling but routes the value stream to the actor collapses identically (hit rate $0.000$, accuracy $0.420$). The effect is categorical --- a $\sim$$79$-point hit-rate gap from a single structural edit --- and establishes two jointly necessary conditions: the loop must be closed through a shared substrate, and the image-grounded stream must command the policy. We trace the failure to credit assignment: a critic that evaluates a different state than the actor occupies cannot teach it. We relate the working loop to the cortico-basal-ganglia--collicular circuit, presenting the correspondence as conceptual motivation rather than literal isomorphism, and use the trained model to generate falsifiable neurophysiological predictions.
\end{abstract}
